% Anti-pattern: pointing at a regression-table row by drawing a TikZ box
% around it. The box adds ink without removing any of the competing
% information. Compare with highlight-good.tex. The table fragment is
% generated by highlight-tables.R.
\documentclass[aspectratio=169]{beamer}
\usetheme{default}
\setbeamertemplate{navigation symbols}{}

\definecolor{accent}{HTML}{107895}
\usecolortheme[named=accent]{structure}

\usepackage[T1]{fontenc}
\usepackage{lmodern}
\usepackage{tabularray}
\usepackage{tikz}
\NewTableCommand{\tinytableDefineColor}[3]{\definecolor{#1}{#2}{#3}}

\begin{document}

\begin{frame}{Results}
  \small
  \begin{table}
\centering
\begin{talltblr}[         %% tabularray outer open
entry=none,label=none,
note{}={+ p < 0.1, * p < 0.05, ** p < 0.01, *** p < 0.001},
]                     %% tabularray outer close
{                     %% tabularray inner open
colspec={Q[]Q[]Q[]Q[]},
hline{2}={1-4}{solid, black, 0.05em},
hline{12}={1-4}{solid, black, 0.05em},
hline{1}={1-4}{solid, black, 0.08em},
hline{14}={1-4}{solid, black, 0.08em},
column{2-4}={}{halign=c},
column{1}={}{halign=l},
}                     %% tabularray inner close
& (1) & (2) & (3) \\
(Intercept) & 79.610*** & 84.080*** & 62.101*** \\
& (2.104) & (5.782) & (9.605) \\
Education & -0.862*** & -0.963*** & -0.980*** \\
& (0.145) & (0.189) & (0.148) \\
Agriculture &  & -0.066 & -0.155* \\
&  & (0.080) & (0.068) \\
Catholic &  &  & 0.125*** \\
&  &  & (0.029) \\
Infant.Mortality &  &  & 1.078** \\
&  &  & (0.382) \\
Num.Obs. & 47 & 47 & 47 \\
R2 & 0.441 & 0.449 & 0.699 \\
\end{talltblr}
\end{table}

  % The box around the Education row IS the anti-pattern shown on this slide.
  % Coordinates are measured from the page center (frame is 160mm x 90mm).
  \begin{tikzpicture}[remember picture, overlay]
    \draw[red, very thick, rounded corners=2pt]
      ([shift={(-4.35cm,1.85cm)}]current page.center) rectangle
      ([shift={(4.40cm,0.65cm)}]current page.center);
  \end{tikzpicture}
\end{frame}

\end{document}
